% GENERATED FILE. Edit canonical lessons, then run npm run generate:books.
% canonical-source-hash: fnv1a64:1c5eeafeb2983abf
% canonical-lessons: BN-C179-chakha, BN-C179-gela, BN-C179-chibano, BN-C179-takano, BN-C179-bachha

\chapter{To taste, To swallow, To chew, To look, To pick}
\label{ch:bn-to-taste-to-swallow-to-chew-to-look-to-pick}
\hlchaptermodality{179}

\begin{chapteropening}
\textbf{By the end of this chapter:} \emph{I can say to taste, to swallow, to chew, to look and to pick.}

Complete the last lesson of chapter 179: I can say to taste, to swallow, to chew, to look and to pick.
\end{chapteropening}

\section[\texorpdfstring{\bn{চাখা} (chākhā)}{চাখা (chākhā)}]{\bn{চাখা} (chākhā) — to taste}

\label{lesson:BN-C179-chakha}



\begin{quote}
Before the new one: say the Bengali for a seller, then the Bengali for a buyer.
\end{quote}

\subsection*{You'll want to know: \bn{চাখা}}
\textbf{\bn{চাখা}} — \emph{chākhā} — \textquotedblleft{}to taste\textquotedblright{}.

\begin{grammarlens}[title={What you've built}]
One more word for this chapter.
\end{grammarlens}

\subsection*{Your turn}
\begin{itemize}
\raggedright
  \item \emph{Say it:} \emph{chākhā}
  \item \emph{Say it:} \emph{chākhā}, once more
  \item \emph{Say it:} say \emph{chākhā}
  \item \emph{Recall:} say the Bengali for a seller, then the Bengali for a buyer, then say \emph{chākhā} again
\end{itemize}

\begin{tcolorbox}[breakable,colback=teal!4,colframe=teal!35!black,title={Before you move on}]
What does \bn{চাখা} mean? (\textquotedblleft{}To taste\textquotedblright{}.) Say it once more.
\end{tcolorbox}
\section[\texorpdfstring{\bn{গেলা} (gelā)}{গেলা (gelā)}]{\bn{গেলা} (gelā) — to swallow}

\label{lesson:BN-C179-gela}



\begin{quote}
Before the new one: say the Bengali for a buyer, then the Bengali for to taste.
\end{quote}

\subsection*{You'll want to know: \bn{গেলা}}
\textbf{\bn{গেলা}} — \emph{gelā} — \textquotedblleft{}to swallow\textquotedblright{}.

\begin{grammarlens}[title={What you've built}]
One more word for this chapter.
\end{grammarlens}

\subsection*{Your turn}
\begin{itemize}
\raggedright
  \item \emph{Say it:} \emph{gelā}
  \item \emph{Say it:} \emph{gelā}, once more
  \item \emph{Say it:} say \emph{gelā}
  \item \emph{Recall:} say the Bengali for a buyer, then the Bengali for to taste, then say \emph{gelā} again
\end{itemize}

\begin{tcolorbox}[breakable,colback=teal!4,colframe=teal!35!black,title={Before you move on}]
What does \bn{গেলা} mean? (\textquotedblleft{}To swallow\textquotedblright{}.) Say it once more.
\end{tcolorbox}
\section[\texorpdfstring{\bn{চিবানো} (chibāno)}{চিবানো (chibāno)}]{\bn{চিবানো} (chibāno) — to chew}

\label{lesson:BN-C179-chibano}



\begin{quote}
Before the new one: say the Bengali for to taste, then the Bengali for to swallow.
\end{quote}

\subsection*{You'll want to know: \bn{চিবানো}}
\textbf{\bn{চিবানো}} — \emph{chibāno} — \textquotedblleft{}to chew\textquotedblright{}.

\begin{grammarlens}[title={What you've built}]
One more word for this chapter.
\end{grammarlens}

\subsection*{Your turn}
\begin{itemize}
\raggedright
  \item \emph{Say it:} \emph{chibāno}
  \item \emph{Say it:} \emph{chibāno}, once more
  \item \emph{Say it:} say \emph{chibāno}
  \item \emph{Recall:} say the Bengali for to taste, then the Bengali for to swallow, then say \emph{chibāno} again
\end{itemize}

\begin{tcolorbox}[breakable,colback=teal!4,colframe=teal!35!black,title={Before you move on}]
What does \bn{চিবানো} mean? (\textquotedblleft{}To chew\textquotedblright{}.) Say it once more.
\end{tcolorbox}
\section[\texorpdfstring{\bn{তাকানো} (tākāno)}{তাকানো (tākāno)}]{\bn{তাকানো} (tākāno) — to look (at)}

\label{lesson:BN-C179-takano}



\begin{quote}
Before the new one: say the Bengali for to swallow, then the Bengali for to chew.
\end{quote}

\subsection*{You'll want to know: \bn{তাকানো}}
\textbf{\bn{তাকানো}} — \emph{tākāno} — \textquotedblleft{}to look (at)\textquotedblright{}.

\begin{grammarlens}[title={What you've built}]
One more word for this chapter.
\end{grammarlens}

\subsection*{Your turn}
\begin{itemize}
\raggedright
  \item \emph{Say it:} \emph{tākāno}
  \item \emph{Say it:} \emph{tākāno}, once more
  \item \emph{Say it:} say \emph{tākāno}
  \item \emph{Recall:} say the Bengali for to swallow, then the Bengali for to chew, then say \emph{tākāno} again
\end{itemize}

\begin{tcolorbox}[breakable,colback=teal!4,colframe=teal!35!black,title={Before you move on}]
What does \bn{তাকানো} mean? (\textquotedblleft{}To look (at)\textquotedblright{}.) Say it once more.
\end{tcolorbox}
\section[\texorpdfstring{\bn{বাছা} (bāchhā)}{বাছা (bāchhā)}]{\bn{বাছা} (bāchhā) — to pick, to choose}

\label{lesson:BN-C179-bachha}



\begin{quote}
Before the new one: say the Bengali for to chew, then the Bengali for to look.
\end{quote}

\subsection*{You'll want to know: \bn{বাছা}}
\textbf{\bn{বাছা}} — \emph{bāchhā} — \textquotedblleft{}to pick, to choose\textquotedblright{}.

\begin{grammarlens}[title={What you've built}]
One more word for this chapter.
\end{grammarlens}

\subsection*{Your turn}
\begin{itemize}
\raggedright
  \item \emph{Say it:} \emph{bāchhā}
  \item \emph{Say it:} \emph{bāchhā}, once more
  \item \emph{Say it:} say \emph{bāchhā}
  \item \emph{Recall:} say the Bengali for to chew, then the Bengali for to look, then say \emph{bāchhā} again
\end{itemize}

\begin{tcolorbox}[breakable,colback=teal!4,colframe=teal!35!black,title={Before you move on}]
What does \bn{বাছা} mean? (\textquotedblleft{}To pick, to choose\textquotedblright{}.) Say it once more.
\end{tcolorbox}
